% Gemini theme
% https://github.com/anishathalye/gemini

\documentclass[final]{beamer}

% ====================
% Packages
% ====================

\usepackage[T1]{fontenc}
\usepackage{lmodern}
\usepackage[size=custom,width=120,height=72,scale=1.0]{beamerposter}
\usetheme{gemini}
\usecolortheme{uw}
\usepackage{graphicx}
\usepackage{booktabs}
\usepackage{tikz}
\usepackage{pgfplots}
\pgfplotsset{compat=1.14}
\usepackage{anyfontsize}

\setmonofont{DejaVu Sans Mono}[Scale=MatchLowercase]
\usepackage{fancyvrb}
% \fvset{fontsize=\small}

\usepackage{tikz}
\usetikzlibrary{cd,decorations.markings,arrows.meta,calc}

\usepackage{
  amsmath,
  amssymb,
  amsthm,
  amsfonts,
  mathtools,
  bbm,
  mathdots,
  bm,        % bold Greek letters
  mathrsfs,  % mathscr
  microtype, % improved formatting
  thmtools,  % fixes theorem referencing
}

\usepackage[normalem]{ulem}

\usepackage{caption}
\usepackage[normalsize]{subfigure}

\usepackage[msc-links, nobysame, lite, url]{amsrefs}

\usepackage[capitalize]{cleveref}
\crefformat{section}{\S{#2#1#3}}
\Crefformat{section}{\S{#2#1#3}}
\crefrangeformat{section}{\S\S{{#3#1#4}} to #5#2#6}
\Crefrangeformat{section}{\S\S{{#3#1#4}} to #5#2#6}
\crefmultiformat{section}{\S{\S{#2#1#3}}}{ and~#2#1#3}{, #2#1#3}{, and~#2#1#3}
\Crefmultiformat{section}{\S{\S{#2#1#3}}}{ and~#2#1#3}{, #2#1#3}{, and~#2#1#3}
\crefformat{equation}{(#2#1#3)}
\Crefformat{equation}{(#2#1#3)}
\crefrangeformat{equation}{(#3#1#4) to (#5#2#6)}
\Crefrangeformat{equation}{(#3#1#4) to (#5#2#6)}
\crefmultiformat{equation}{(#2#1#3)}{ and~(#2#1#3)}{, (#2#1#3)}{, and~(#2#1#3)}
\Crefmultiformat{equation}{(#2#1#3)}{ and~(#2#1#3)}{, (#2#1#3)}{, and~(#2#1#3)}

\usepackage[capitalize]{cleveref} 
% \mathtoolsset{showonlyrefs} % this is not compatible with cleveref; instead:
\usepackage{autonum}

% \renewcommand{\arraystretch}{1.6}
% \makeatletter
% \renewcommand*\env@matrix[1][*\c@MaxMatrixCols c]{
%   \hskip -\arraycolsep
%   \let\@ifnextchar\new@ifnextchar
%   \array{#1}}
% \makeatother

\usepackage{style}
\newcommand*{\defn}[1]{\textbf{#1}}

% ====================
% Lengths
% ====================

% If you have N columns, choose \sepwidth and \colwidth such that
% (N+1)*\sepwidth + N*\colwidth = \paperwidth
\newlength{\sepwidth}
\newlength{\colwidth}
\setlength{\sepwidth}{0.025\paperwidth}
\setlength{\colwidth}{0.3\paperwidth}

\newcommand{\separatorcolumn}{\begin{column}{\sepwidth}\end{column}}

% ====================
% Title
% ====================

\title{Formalizing polyhedral geometry in Lean}
\author{Seven Lewis \and George Peykanu \and Caelan Ritter \and Freda Zhang}
\institute[]{University of Washington}

% ====================
% Footer (optional)
% ====================

\footercontent{Washington Experimental Mathematics Lab, Spring 2025}
% (can be left out to remove footer)

% ====================
% Logo (optional)
% ====================

% use this to include logos on the left and/or right side of the header:
% \logoright{\includegraphics[height=7cm]{logo1.pdf}}
% \logoleft{\includegraphics[height=7cm]{logo2.pdf}}

% ====================
% Body
% ====================

\begin{document}

\addtobeamertemplate{headline}{}{\begin{tikzpicture}[remember picture, overlay]
  \node [anchor=north east, inner sep=1cm, xshift=-0.5cm] at (current page.north east)
  {\includegraphics[height=6.8cm]{figures/qrcode.png}};
\end{tikzpicture}}

\begin{frame}[fragile]
    \begin{columns}[t]
        \separatorcolumn

        \begin{column}{\colwidth}

            \begin{block}{Background}    
                The primary focus of this quarter was in the formalization of two important theorems in polyhedral combinatorics: Carath\'eordory's Theorem, and the Hyperplane Separation Theorem. The formalization of each was successfully accomplished in the span of the ten-week cycle, as well as preliminary work in the formalization of Polar Spaces and Dual Topologies, with the enduring focus of the project, to prove the ``main'' theorem of Polyhedral Geometry, that is, that a subset of a vector space is a polyhedral cone if and only if it is the conical hull of a finite set of vectors.

                This result provides a combinatorial approach to the finite generation of polyhedra, allowing one to choose a subset of points from a vector space and in doing so, to uniquely determine a polyhedral cone.
            \end{block}

            \begin{block}{Preliminaries}
                Have $V = \RR^d$ be the ambient finite-dimensional vector space. The \textbf{halfspace} is commonly defined as the set of points on one side of an affine hyperplane:
                $$ \{\mathbf{x} \in V \;|\; \mathbf{ax} \leq b\}$$
                where $\mathbf{a} \in \RR^d$ is the normal vector of the hyperplane with $b \in \RR$ as the affine constant term.

                We define the \textbf{finite polyhedron} as the intersection of finitely many halfspaces $\{H_1,H_2,\dots,H_n\}$:
                $$ \bigcap_{1 \leq i \leq n} H_i. $$
                A \textbf{convex} set $S \subseteq V$ is any set with the property that given two elements $\mathbf{x}, \mathbf{y} \in S$, we have
                $$ (1 - t)\mathbf{x} + t\mathbf{y} \in S \quad \text{ for all real } t \in [0,1].$$
                A \textbf{convex hull} is the closure of any $S \subseteq V$ under the property of convexity. Equivalently, $\text{conv}\,S$ is the smallest convex set of $V$ which contains $S$. A \textbf{convex combination} of elements $x_1,x_2,\dots,x_n \in V$ is a linear combination where the coefficients $\lambda_i$ are nonnegative and sum to 1. 

                A \textbf{polytope} can be defined as a bounded polyhedron or, equivalently, the convex hull of finitely many points in $\RR^d$. One eventual goal of our project is to prove this equivalency. 
                
                It turns out that conical sets are a nice generalization of the idea of a \textbf{convex set}. As a result, much of our formulation relies on the definitions of the following: A \textbf{conical combination} of elements $x_1,x_2,\dots,x_n \in V$ is a linear combination where the coefficients $\lambda_i$ are nonnegative. A \textbf{conical set} is simply closed under the operation of conical combination. A \textbf{conical hull} is the closer of the set under finite conical combinations.
                 
                In proving the Hyperplane Separation Theorem, two widely-known theorems were crucial: the \textbf{ Heine-Borel Theorem}, and the existence of an infimum on a compact set.
            \end{block}

            \begin{block}{Definitions translated to Lean}
            \vspace{-0.5em}
\begin{Verbatim}
def Halfspace : Prop := ∃ (f : V →ₗ[ℝ] ℝ) (c : ℝ), s = { x | f x ≤ c }
def Polyhedron : Prop := ∃ (I : Type*) (H : I → Set V), Finite I ∧ 
  (∀ i : I, Halfspace (H i)) ∧ s = ⋂ (i : I), H i
def Cone : Prop := s.Nonempty ∧ ∀ c ≥ (0 : ℝ), ∀ x ∈ s, c • x ∈ s
def conicalHull.{u} : Set V := { x | isConicalCombo.{_, u} s x }
def isConicalCombo (x : V) : Prop :=  ∃ (ι : Type*) (t : Finset ι) 
  (a : ι → ℝ) (v : ι → V), isConicalCombo' s x t a v
def isConicalCombo' (x : V) {ι : Type*} (t : Finset ι) (a : ι → ℝ) (v : ι → V)
  : Prop := (∀ i ∈ t, a i = 0 ∨ 0 ≤ a i ∧ v i ∈ s) ∧ x = ∑ i ∈ t, a i • v i

\end{Verbatim}
        \end{block}
           
            
        \end{column}

        \separatorcolumn

        \begin{column}{\colwidth}
          \vspace{1em}
          \begin{verbatim}
theorem Metric.isCompact_iff_isClosed_bounded.{u} [...] :
  IsCompact s ↔ IsClosed s ∧ Bornology.IsBounded s
theorem IsCompact.exists_isMinOn : {s : Set β} (hs : IsCompact s) (ne_s :
  s.Nonempty) {f : β → α}(hf : ContinuousOn f s) : ∃ x ∈ s, IsMinOn f s x
\end{verbatim}

\vspace{1em}


            \begin{exampleblock}{A sample proof}
            \vspace{-0.5em}
\begin{Verbatim}
  theorem poly_convex : ∀ (s : Set V), Polyhedron s → Convex ℝ s := by
    intro s h_s_poly
    unfold Polyhedron at h_s_poly
    rcases h_s_poly with ⟨I, H, h_I_finite, h_Hi_halfspace, rfl⟩
    apply convex_iInter
    intro i
    exact halfspace_convex _ (h_Hi_halfspace i)
\end{Verbatim}
            \end{exampleblock}


            \begin{alertblock}{Carath\'eodory's Theorem (1.3.2)}
              Let $S \subseteq V = \RR^d$. For all $x \in \text{cone}\,S$, $x$ is a conical combination of at most $d$ elements of $S$. For all $x \in \text{conv}\,S$, $x$ is a convex combination of at most $d + 1$ elements of $S$. 
            \end{alertblock}

            \begin{figure}[htbp]
                    \centering
                    \begin{tikzpicture}[scale=1.5, thick, >=Stealth, decoration={markings,mark=at position 0.5 with {\arrow[xshift={4pt + 2.25\pgflinewidth}]{Latex[scale=1.5]},}}]
  \clip (-2.25,-2.25) rectangle (2.25, 2.25);

  % 1) The six true vertices of a regular hexagon of radius 1
  \foreach \i in {0,...,5} {
    \coordinate (P\i) at ({2*cos(60*\i)},{2*sin(60*\i)});
  }

  % 2) Draw each extended edge in dotted style and name it E<i>
  \foreach \i/\j in {0/1,1/2,2/3,3/4,4/5,5/0} {
    \draw[name path=E\i,dotted]
    ($(P\i)!-2!(P\j)$) -- ($(P\j)!-2!(P\i)$);
  }

  % 4) Shade the triangle P0–P2–P4
  \fill[lightpurple] (P0) -- (P1) -- (P4) -- cycle;

  % 5) Construct two medians of that triangle and name their intersection G
  \coordinate (m02) at ($ (P0)!0.5!(P2) $);   % midpoint of P0–P2
  \coordinate (m24) at ($ (P2)!0.5!(P4) $);   % midpoint of P2–P4

  \path[name path=med1] (P4) -- (m02);  
  \path[name path=med2] (P0) -- (m24);

  \path[name intersections={of=med1 and med2, by=G}];

  % 6) Draw a dot at the centroid G
  \coordinate (Gshift) at,(,(G)+(-0.3cm,0.15cm),);\quad
  \fill (Gshift) circle (1pt);

  % 3) Over‐draw the actual hexagon edges in solid
  \draw[very thick] (P0)--(P1)--(P2)--(P3)--(P4)--(P5)--cycle;

  % (Optional) put dots at the true hexagon vertices, too
  \foreach \i in {0,...,5} {
    \fill (P\i) circle (1pt);
  }
\end{tikzpicture}

%%% Local Variables:
%%% mode: latex
%%% TeX-master: "../poster"
%%% End:

                    \caption{An arbitrary point in the convex hull of 6 points is the convex combination of 3.\label{fig:caratheodory}}
            \end{figure}

            \begin{block}{Formalizing Carath\'eodory's Theorem}
              An important introductory result of conical sets is the notion that any particular vector in the hull of some set $S$ can be constructed as the conical combination of a smaller subset of $S$. This idea can then be extended to the convex case with one additional vector.

              Our approach to formulating this into Lean is as follows:
\begin{Verbatim}
    theorem caratheordory (s : Set V) : ∀ x ∈ conicalHull.{_,0} s, 
      isConicalCombo_aux s x (finrank ℝ V) 
\end{Verbatim}
              \verb|isConicalCombo_aux| is our definition of the previously defined \verb|isConicalCombo| that exposes the number of elements needed in the conical combination of \verb|x|. Since \verb|finrank ℝ V| is simply the dimension of the space, it corresponds directly to $d$ in the original theorem.

              The proof of this in Lean involves a variety of mathematical tools. We have \verb|v| be the set of vectors involved in the conical combination of \verb|x| (which is a subset of \verb|s|). First, we induct on minimum number of elements \verb|N| in \verb|s| required in a conical combination to construct \verb|x|. Then, we do case work on the number of elements \verb|N+1| in \verb|s|: if \verb|N+1| is less than \verb|finrank ℝ V|, we are done; otherwise, we can show that there exists linear dependence in \verb|v| and work to remove the dependent vector while preserving nonnegativity of the coefficients of the remaining vectors in \verb|v|.  
            \end{block}    
            
            

        \end{column}

        \separatorcolumn

        \begin{column}{\colwidth}

            
          \begin{alertblock}{Hyperplane Separation Theorem}
            Let $A$, and $B$ be two convex, closed subsets of $V$, with $A$ compact. Then $
                \text{there exists } \; c\in \mathbb{R}, f : V\to \mathbb{R} \; \text{such that} \; f(a) < c \; \text{for all}\; a\in A, \ c < f(b)\; \text{for all}\; b\in B
              $ Where $f$ is a linear functional from $V\to \mathbb{R}$.
        \end{alertblock}

        \begin{figure}[htbp]
            \centering
            \begin{tikzpicture}[scale=3, very thick, >=Stealth, decoration={markings,mark=at position 0.5 with {\arrow[xshift={4pt + 2.25\pgflinewidth}]{Latex[scale=1.5]},}}]
  \clip (-0.25,-0.25) rectangle (4,2);
  \filldraw[fill=lightpurple] (0,0) circle [radius=1];
  \node[anchor=south west] at (0,0) {$A$};
  \filldraw[fill=lightpurple] (5,0) -- (3,0.5) -- (45:2) -- (0, 3) -- (5,3) -- cycle;
  \node[anchor=north east] at (3.75,1.75) {$B$};
  \filldraw (45:1) circle [radius=0.04];
  \filldraw (45:2) circle [radius=0.04];
  \draw[dotted] (45:1) -- (45:2);
  \draw (45:1.5) + (-45:3) -- + (135:3);
  % \filldraw (0,-2) circle [radius=0.06];
  % \draw[postaction={decorate}] (-90:2) arc [start angle=-90, end angle=90, radius=2];
  % \draw[postaction={decorate}] (-90:2) arc [start angle=-90, end angle=-270, radius=2];
  % \draw[postaction={decorate}] (0,2) -- (0,-2);
  % \draw (0,2.6) node {$q$};
  % \draw (180:2.6) node {$e_{1}$};
  % \draw (0:2.6) node {$e_{2}$};
  % \draw (0.6,0) node {$e_{3}$};
\end{tikzpicture}

%%% Local Variables:
%%% mode: latex
%%% TeX-master: "../poster"
%%% End:

            \caption{A separating hyperplane determined by infimal points. \label{fig:hyperplane-separation}}
        \end{figure}

\begin{block}{Formal Statement in Lean4}
\vspace{-0.5em}
\begin{Verbatim}
theorem hyperplane_separation  (A B : Set V) (hA : Convex ℝ A)(hB : Convex 
ℝ B) (hclosed : IsClosed A ∧ IsClosed B )(hNempty : A.Nonempty ∧ B.Nonempty) 
(hA_Bounded: IsBounded A) (hAB : Disjoint A B) : ∃ (f : V →ₗ[ℝ] ℝ) (c : ℝ),
(∀ a ∈ A, f a < c) ∧ (  ∀ b ∈ B, c < f b)
\end{Verbatim}
\end{block}

        \begin{block}{Approaching Hyperplane Separation Theorem}
            The Hyperplane Separation Theorem is a key result in proving the existence of certain polyhedral constructs.
            \\\\The formal proof involves exhibiting an infimal distance between the sets, and the corresponding points, $a_{\text{min}}, b_{\text{min}} \in V$ which define this infimal distance.
            $f:= b_{\text{min}}-a_{\text{min}}$ is thus a vector in $V$, and defining a functional by $\langle f, \_\rangle$ constructs a linear functional from $V\to \mathbb{R}$. Any hyperplane to which $f$
            is normal, and which does not intersect $A$ or $B$, then separates the two sets as desired, after unraveling the convexity of each set.
            \\\\Convexity is used to provide contradiction, by providing a linear combination of points within each set which is simultaneously the norm of a vector,
            and negative for small coefficients. In the process of formalization, this means providing an explicit value for which the contradiction arises.         
        \end{block}
        
         
        \begin{exampleblock}{Corollary. Farkas's Lemma}
          Let $C$ be a closed, convex cone, in $V$, $u\in V\setminus C$. Then there exists a separating hyperplane between $u$ and $C$, defined by a linear functional $y$ in the space of functions $V\to \mathbb{R}$, such that $
          y(u)> 0, \; \text{for all } x\in C, \; \text{and} \; y(x) \le 0
          $
        \end{exampleblock}
            This result is  a combination of the Hyperplane Separation Theorem, for which the conclusion ``There exists $
              c\in \mathbb{R}, y \in V^* \text{ such that } y(a) \le c\; \text{for all } x\in C, \; c < f(x)\; \text{for all }  b\in B
            $'' immediately follows, and a translation lemma for cones proven earlier in the project.
\begin{Verbatim}
lemma farkas (u: V)(C: Set V) (convC: Convex ℝ C) (closedC: IsClosed C)
(coneC: Cone C): u ∉ C → ∃(y : V →ₗ[ℝ] ℝ), (y u > 0) ∧ (∀ x ∈ C, y x ≤ 0)
\end{Verbatim}  
           
            
        \begin{block}{References}
        \nocite{liunotes}
        \begin{bibdiv}
        \begin{biblist}
          \bib{liunotes}{webpage}{
              title={Polyhedral combinatorics},
              author={Liu, Gaku},
              accessdate={3 June 2023},
              url={https://sites.google.com/view/gakuliu/home}
            }
        \end{biblist}
        \end{bibdiv}
        \end{block}
\end{column}

\separatorcolumn
\end{columns}
\end{frame}
\end{document}

%%% Local Variables:
%%% mode: LaTeX
%%% TeX-master: t
%%% End:
